\section{Methodology}

The CMatrix framework is designed to transform high-level security objectives into safe, executable, and self-correcting attack chains. This section details our formal threat model, system architecture, and the composite reasoning orchestration algorithm.

\subsection{Threat Model and Security Boundaries}
We define a formal threat model where the adversary is an AI-augmented red team agent tasked with vulnerability discovery. As illustrated in Fig.~\ref{fig:threat-model}, we establish a \textbf{Trust Boundary} around the Orchestrator and the HITL Safety Gate. We assume the agent has network-level access to the target infrastructure but is constrained by a risk-aware policy $P$ that intercepts any command with a disruption score above a defined threshold.

% ============================================================
% FIGURE 01: FORMAL THREAT MODEL
% File: figure-01-threat-model.tex
% Section: §3.1 Threat Model
% Description: Visual representation of the adversary, attack surface,
%              target system, and HITL safety boundaries.
% ============================================================

\begin{figure*}[htbp]
  \centering
  \begin{tikzpicture}[
    node distance=2cm,
    font=\small,
    box/.style={rectangle, draw, rounded corners, minimum width=2.5cm, minimum height=1cm, align=center, fill=gray!5},
    agent/.style={box, fill=blue!10, thick},
    gate/.style={diamond, draw, fill=red!10, minimum width=2.5cm, minimum height=1.5cm, align=center},
    target/.style={box, fill=green!10, thick},
    boundary/.style={draw, dashed, gray!60, inner sep=0.5cm}
  ]

    % Nodes
    \node (user) [box] {Security Operator \\ (Human-in-the-Loop)};
    \node (orch) [agent, below of=user, yshift=-0.5cm] {\textbf{CMatrix Orchestrator} \\ (Reasoning Suite)};
    \node (gate) [gate, right of=orch, xshift=2.5cm] {HITL \\ Safety Gate};
    \node (tools) [box, below of=gate, yshift=-0.5cm] {Security Tool Registry \\ (Nmap, Metasploit, etc.)};
    \node (target) [target, right of=gate, xshift=2.5cm] {Target Infrastructure \\ (Enterprise Network)};

    % Boundaries
    \begin{scope}[on background layer]
      \node [boundary, fit=(user) (orch) (gate) (tools), label=above:{\textbf{Trust Boundary}}] (trust) {};
    \end{scope}

    % Flows
    \draw[->, thick] (user) -- node[right] {Objective} (orch);
    \draw[->, thick] (orch) -- node[above, align=center] {Proposed \\ Actions} (gate);
    \draw[->, thick] (gate) -- node[left, align=center] {Approved \\ Tools} (tools);
    \draw[->, thick, red!60] (gate) -- node[above] {Veto} (user);
    \draw[->, thick] (tools) -- node[above, align=center] {Exploitation \\ Probes} (target);
    \draw[->, dashed] (target) -- node[below] {Observations} (orch);
    \draw[->, dashed] (orch) -- node[above right, align=center] {Risk Telemetry} (user);

    % Legend
    \node[anchor=north west] at (current bounding box.south west) [yshift=-0.5cm] {
      \begin{tabular}{ll}
        \tikz \draw[->, thick] (0,0) -- (0.5,0); & Control Flow \\
        \tikz \draw[->, dashed] (0,0) -- (0.5,0); & Data/Observation Flow \\
      \end{tabular}
    };

  \end{tikzpicture}
  \caption{Formal threat model of CMatrix. The security operator defines objectives for the autonomous orchestrator, while the HITL safety gate provides a hard trust boundary between the reasoning suite and the target infrastructure.}
  \label{fig:threat-model}
\end{figure*}


\subsection{System Architecture}
The technical architecture of CMatrix follows a "Master-Worker" paradigm implemented using LangGraph's stateful state machine. As shown in Fig.~\ref{fig:system-architecture}, the \textbf{Orchestrator} coordinates specialized reasoning services and delegates operational tasks to a registry of worker agents via a Celery background task queue. State persistence is maintained in Redis for ephemeral reasoning traces and PostgreSQL for permanent audit logs.

% ============================================================
% FIGURE 02: SYSTEM ARCHITECTURE
% File: figure-02-system-architecture.tex
% Section: §3.2 System Architecture
% Description: Technical architecture of the CMatrix execution spine.
% ============================================================

\begin{figure}[htbp]
  \centering
  \begin{tikzpicture}[
    node distance=1.5cm,
    font=\footnotesize,
    box/.style={rectangle, draw, rounded corners, minimum width=2cm, minimum height=0.8cm, align=center},
    service/.style={box, fill=blue!5},
    storage/.style={cylinder, draw, shape border rotate=90, minimum width=1.5cm, minimum height=1.2cm, fill=yellow!5},
    logic/.style={box, fill=orange!10, thick},
    worker/.style={box, fill=green!5, dashed}
  ]

    % Frontend/API
    \node (api) [service] {API Layer \\ (FastAPI)};

    % Orchestrator
    \node (orch) [logic, below of=api, yshift=-0.5cm] {\textbf{Orchestrator} \\ (LangGraph StateGraph)};

    % Reasoning Services
    \node (tot) [service, left of=orch, xshift=-2.5cm] {ToT \\ Strategy};
    \node (rewoo) [service, below of=tot] {ReWOO \\ Planner};
    \node (refl) [service, below of=rewoo] {Reflexion \\ Engine};

    % State
    \node (redis) [storage, right of=orch, xshift=2.5cm] {State Cache \\ (Redis)};
    \node (db) [storage, below of=redis, yshift=-0.5cm] {Audit Log \\ (PostgreSQL)};

    % Workers
    \node (celery) [box, below of=orch, yshift=-1cm] {Task Queue \\ (Celery)};
    \node (agents) [worker, below of=celery] {Specialized Agents \\ (Network, Web, API, etc.)};

    % Flows
    \draw[<->, thick] (api) -- (orch);
    \draw[->] (orch) -- (tot);
    \draw[->] (orch) -- (rewoo);
    \draw[->] (orch) -- (refl);
    \draw[<->] (orch) -- (redis);
    \draw[->] (orch) -- (db);
    \draw[->, thick] (orch) -- (celery);
    \draw[->, thick] (celery) -- (agents);
    \draw[->, dashed] (agents.east) .. controls +(2,-1) and +(2,1) .. (orch.east) node[midway, right] {Observations};

  \end{tikzpicture}
  \caption{CMatrix System Architecture. The Orchestrator manages the state machine, delegating complex reasoning to specialized services and executing long-running security tasks via a background task queue.}
  \label{fig:system-architecture}
\end{figure}


\subsection{Composite Reasoning Pipeline}
The core of our methodology is the sequential composition of three advanced reasoning patterns, as detailed in Alg.~\ref{alg:composite-reasoning}. 

% ============================================================
% ALGORITHM 01: COMPOSITE REASONING PIPELINE (REVISED FOR COMPATIBILITY)
% File: algorithm-01-composite-reasoning.tex
% Description: Formal flow of the ToT + ReWOO + Reflexion composition.
% ============================================================

\begin{figure}[htbp]
\centering
\begin{minipage}{0.48\textwidth}
\hrule
\vspace{2pt}
\textbf{Algorithm 1} Composite Reasoning Orchestration
\vspace{2pt}
\hrule
\vspace{5pt}
\small
\begin{enumerate}[label=\arabic*:, leftmargin=2em]
    \item \textbf{Input:} Objective $O$, Target Context $C$, Risk Policy $P$
    \item \textbf{Output:} Attack Report $R$, Audit Logs $L$
    \item $S \gets \text{TreeOfThoughts}(O, C)$ \hfill \textit{// Select optimal strategy}
    \item $B \gets \text{ReWOOPlanner}(O, S, C)$ \hfill \textit{// Generate blueprint}
    \item \textbf{while} $B$ is not empty \textbf{do}
    \item \quad $A \gets B.\text{next\_action}()$
    \item \quad \textbf{if} $\text{RiskAssessment}(A, P) > \text{Threshold}$ \textbf{then}
    \item \quad \quad $A \gets \text{HumanInTheLoopApproval}(A)$
    \item \quad \textbf{end if}
    \item \quad $Obs \gets \text{ExecuteAgentTask}(A)$
    \item \quad $L.\text{log}(A, Obs)$
    \item \textbf{end while}
    \item $I \gets \text{ReflexionEngine}(L, O)$ \hfill \textit{// Identify gaps}
    \item \textbf{if} $I \neq \emptyset$ \textbf{then}
    \item \quad $B' \gets \text{GenerateCorrectivePlan}(I)$
    \item \quad \textbf{goto} Line 5 with $B = B'$
    \item \textbf{end if}
    \item $R \gets \text{SynthesizeReport}(L)$
    \item \textbf{return} $R$
\end{enumerate}
\vspace{2pt}
\hrule
\end{minipage}
\caption{The composite reasoning orchestration pipeline. The process combines proactive planning (ReWOO) with iterative self-correction (Reflexion) and strategic lookahead (ToT).}
\label{alg:composite-reasoning}
\end{figure}


\begin{enumerate}
    \item \textbf{Strategy Selection (ToT)}: The system first performs a Tree of Thoughts search to evaluate multiple attack strategies (e.g., Stealthy vs. Aggressive) against five domain-specific heuristics.
    \item \textbf{Blueprint Planning (ReWOO)}: Based on the selected strategy, the ReWOO planner generates a dependency-aware blueprint. This reduces token overhead by batches of tool calls before execution.
    \item \textbf{Executable Reflection}: After tool execution, the Reflexion engine performs a gap analysis on the observations. Crucially, it converts identified weaknesses into structured \texttt{ImprovementAction} objects that trigger iterative planning updates (Fig.~\ref{fig:reasoning-flow}).
\end{enumerate}

% ============================================================
% FIGURE 03: REASONING FLOW
% File: figure-03-reasoning-flow.tex
% Section: §3.3 Composite Reasoning
% Description: Detailed flow of the Strategy -> Planning -> Execution -> Reflection loop.
% ============================================================

\begin{figure}[htbp]
  \centering
  \begin{tikzpicture}[
    scale=0.85, transform shape,
    node distance=1.8cm,
    font=\small,
    phase/.style={rectangle, draw, rounded corners, minimum width=3cm, minimum height=1cm, align=center, fill=blue!5, thick},
    data/.style={ellipse, draw, fill=yellow!10, minimum width=2cm},
    arrow/.style={->, >=stealth, thick}
  ]

    % Nodes
    \node (obj) [data] {Objective $O$};
    \node (tot) [phase, below of=obj] {\textbf{ToT Strategy Selection} \\ (Search \& Heuristics)};
    \node (rewoo) [phase, below of=tot] {\textbf{ReWOO Planning} \\ (Decoupled Blueprint)};
    \node (exec) [phase, below of=rewoo] {\textbf{Execution \& HITL} \\ (Specialized Agents)};
    \node (refl) [phase, below of=exec] {\textbf{Reflexion Engine} \\ (Gap Analysis)};

    % Flows
    \draw [arrow] (obj) -- (tot);
    \draw [arrow] (tot) -- node[right] {Strategy $S$} (rewoo);
    \draw [arrow] (rewoo) -- node[right] {Blueprint $B$} (exec);
    \draw [arrow] (exec) -- node[right] {Observations $Obs$} (refl);

    % Loop back
    \draw [arrow, red!70] (refl.west) .. controls +(-2,0) and +(-2,0) .. node[left, align=center] {Corrective \\ Actions $I$} (rewoo.west);

    % Exit
    \node (rep) [data, right of=refl, xshift=2.5cm] {Final Report $R$};
    \draw [arrow] (refl) -- node[above] {Success} (rep);

  \end{tikzpicture}
  \caption{Composite Reasoning Flow. The process initiates with ToT strategy selection, followed by ReWOO blueprint generation. The Reflexion engine identifies gaps in observations and triggers iterative planning updates until the objective is met.}
  \label{fig:reasoning-flow}
\end{figure}

